\documentclass[book.tex]{subfiles}

\begin{document}

\chapter{Scattering Theory}

\label{chap:quantum_harmonic}
\noindent\rule{11cm}{0.4pt} \\
\textbf{Prerequisites:}
\autoref{chap:harmonic},
\autoref{chap:quantum_mechanics}, \autoref{chap:position_and_momentum} \\
\textbf{Difficulty Level:} ** \\
\noindent\rule{11cm}{0.4pt}
\vspace{5mm}
%\textcolor{red}{TODO: Cite Sakurai}

A scattering process studies the collision of particles. Mathematically for a scattering process, we start from an input state of particles and end with a final state. The scattering process can be modeled either classically or quantum mechanically.


\section{Classical Scattering}

Broadly, classical scattering can be broken into elastic and ineleastic scattering. Elastic scattering conserves energy and treats particles as point objects that don't change in the collision. Inelastic collisions allow for the possibility of deformation and the loss of energy during the collision process.

\section{Quantum Mechanical Scattering}

The transformation process can be modeled as a time dependent perturbation since the shifting event can often be tightly focused in time (a particle collision for example)

Let us start by writing the Hamiltonian to model the scattering process as as
\begin{align}
    H &= H_0 + V_I(r) \\
    H_0 &= \frac{p^2}{2m}
\end{align}
The eigenvalues of the kinetic energy component $H_0$ are
\begin{align}
E_k &= \frac{\hbar^2k^2}{2m}
\end{align}

Let $\alpha(t)$ denote the state at time $t$. We have the transition equation

\begin{align}
    |\alpha(t)\rangle_I &= U_I(t)|\alpha(t_0) \rangle_I\\
\end{align}
Here $U_I$ satisfies the equation %(\textcolor{red}{TODO: Add the derivation of this equation to the perturbation theory part}
\begin{align}
    i\hbar\frac{\partial}{\partial t} U_I(t, t_0) &= V_I(t) U_I(t, t_0)
\end{align}
where
\begin{align}
    U_I(t_0) &= I \\
    V_I(t) &= \exp(iH_0t/\hbar)V\exp(-iH_0t/\hbar)
\end{align}

This equation has solution
\begin{align}
U_I(t) &= 1 - \frac{i}{\hbar}\int_{t_0}^T V_I(t')U_I(t', t_0)dt'
\end{align}
Note that this equation is integrating operators. %(\textcolor{red}{TODO: Is the Lebesgue integral defined over operator space? I think so.}. 
We can then derive the following general solution for the transmission probability from basis state $i$ to basis state $n$%\textcolor{red}{TODO: I don't fully follow the derivation jump Sakurai makes}

\begin{align}
    \langle n | U_I(t)| i \rangle = \delta_{ni} - \frac{i}{\hbar}\sum_m \langle n|V|m \rangle \int_{t_0}^t e^{i\omega_{nm}t'} \langle m | U_I(t')|i \rangle dt'
\end{align}

For specific scattering potentials, we can then derive specific scattering amplitudes. Such scattering amplitudes arise often in high energy physics applications.
\end{document}